%%=============================================================================
%% Inleiding
%%=============================================================================

\chapter{\IfLanguageName{dutch}{Inleiding}{Introduction}}%
\label{ch:inleiding}

%De inleiding moet de lezer net genoeg informatie verschaffen om het onderwerp te begrijpen en in te zien waarom de onderzoeksvraag de moeite waard is om te onderzoeken. In de inleiding ga je literatuurverwijzingen beperken, zodat de tekst vlot leesbaar blijft. Je kan de inleiding verder onderverdelen in secties als dit de tekst verduidelijkt. Zaken die aan bod kunnen komen in de inleiding:

%\begin{itemize}
%  \item context, achtergrond
%  \item afbakenen van het onderwerp
%  \item verantwoording van het onderwerp, methodologie
%  \item probleemstelling
%  \item onderzoeksdoelstelling
%  \item onderzoeksvraag
%  \item \ldots
%\end{itemize}

Mainframe is een van de oudste technologieën uit het computertijdperk. Toch wordt ze nog elke dag door miljoenen mensen gebruikt. De grootste bedrijven en financiële instellingen gebruiken vandaag nog steeds diezelfde technologie. Denk maar aan een grote bank of supermarktketen. Zij vertrouwen allemaal op de stabiele technologie die mainframes bieden. Uiteraard worden de hardware en software voortdurend verbeterd. Zo is er nu bijvoorbeeld al AI en kwantum-safe technologie verbonden aan de mainframe. Ook worden mainframes moderner gemaakt door gebruik te maken van bijvoorbeeld API's (Application Programming Interface) en het incorporeren van UNIX functionaliteiten met de hulp van USS (UNIX System Services). De verbetering en modernisatie van de mainframetechnologie lijkt echter geen weerslag te vinden in de algemene perceptie van mainframes. Jarenlang zijn mainframes beschouwd als een oude technologie die geen toekomst heeft \autocite{Sagers2018}. Daardoor hebben heel weinig mensen zich gespecialiseerd in mainframes. Momenteel is Rexx de meest gebruikte scripting taal op mainframe \autocite{Gilio2021}. Hoewel Rexx ook kan gebruikt worden op andere platformen dan mainframe, is het geschreven om perfect samen te werken met mainframes en wordt het dus ook voornamelijk daarvoor gebruikt. Het grootste aandeel van mensen dat met Rexx kan werken, is reeds tewerkgesteld, en velen gaan binnenkort met pensioen. Het gevolg is dat bedrijven nu kampen met een groot tekort aan beschikbare werkkrachten die overweg kunnen met die unieke technologie \autocites{NgoYe2018}.

\section{\IfLanguageName{dutch}{Probleemstelling}{Problem Statement}}%
\label{sec:probleemstelling}

Zowel de modernisatie als het tekort aan werkkrachten zorgen ervoor dat bedrijven en experts voor een keuze komen te staan. Die keuze is tussen oudere technologieën, zoals Rexx, die stabiliteit bieden en die diep geïntegreerd zijn in het mainframe systeem, en modernere technologieën, zoals Python, die meer flexibiliteit bieden en meer gekend zijn onder nieuwe mainframers.

Een van de mogelijke oplossingen om het tekort aan werkkrachten op te lossen is het aanbieden van meer mainframe opleidingen waar Rexx wordt aangeleerd, zowel via hoger onderwijs, als intern bij de bedrijven zelf. 

Aangezien er momenteel maar heel weinig scholen zijn die Rexx aanleren aan hun studenten, moeten bedrijven voorlopig op zoek naar een andere manier om jong talent toch naar mainframes te lokken \autocite{Blondeel2025}. Een mogelijke andere manier is om de modernisatie te omarmen, en mainframes toegankelijker en aantrekkelijker te maken door gebruik te maken van modernere technologieën die jongere mensen vaker al weten te gebruiken. Daar komt Python naar boven. Python is een alom gekende programmeer – en scripting taal, zelfs bij mensen met een beperkte achtergrond in IT. Het is dus de perfecte taal om zoveel mogelijk mensen te bereiken. Daarbovenop is Python ideaal voor modernisatie, aangezien het de deur opent om de mainframe te laten samenwerken met andere systemen. Bijvoorbeeld volgens \textcite{Gilio2021} is Python nu al een belangrijk deel van de modernisatie op mainframe en zal Python ook een groot deel vormen van de toekomstige modernisatie op mainframe. 

Het is echter niet zo vanzelfsprekend om scripts geschreven in Rexx, die al 40 jaar oud zijn, zomaar te vervangen door Python, een in vergelijking veel modernere taal. Of Python alle Rexx code zomaar kan vervangen is nog onduidelijk, maar hoogst onwaarschijnlijk. Het is dus een lastige beslissing voor bedrijven om ervoor te kiezen de stabiliteit van Rexx zomaar te vervangen door Python, ook al zou dat de nodige modernisatie vergemakkelijken en nieuwe werkkrachten kunnen aantrekken.


\section{\IfLanguageName{dutch}{Onderzoeksvraag}{Research question}}%
\label{sec:onderzoeksvraag}

Dit onderzoek wil bijdragen tot een oplossing voor die problematiek door de volgende onderzoeksvraag te beantwoorden: Hoe verhouden REXX en Python zich voor automatisatie op z/OS vanuit technisch en praktisch oogpunt. 
Om deze vraag te beantwoorden moeten volgende deelvragen beantwoord worden:

\begin{itemize}
    \item Waarom is er nood aan modernisatie van mainframes?
    \item Wat doet automatisatie op mainframes?
    \item Wat zijn de voor- en nadelen van Rexx?
    \item Wat zijn de voor- en nadelen van Python op z/OS?
    \item In hoeverre kan Python nu al Rexx vervangen?
\end{itemize}


\section{\IfLanguageName{dutch}{Onderzoeksdoelstelling}{Research objective}}%
\label{sec:onderzoeksdoelstelling}

De onderzoeksvragen leiden tot de volgende onderzoeksdoelstellingen:

\begin{itemize}
    \item Achterhalen wat de voornaamste redenen zijn voor modernisatie in mainframes.
    \item Begrijpen wat automatisatie doet op mainframes en aantonen hoe Rexx daarin de meeste gebruikte scripting taal is geworden.
    \item Python en Rexx vergelijken op basis van vooropgestelde requirements, zoals complexiteit en moeilijkheid van onderhoud.
    \item Via een vergelijkende studie en PoC's achterhalen of Python nu al een waardige vervanger kan zijn van Rexx, en zo ja, in welke use cases.
\end{itemize}

Het uiteindelijke doel van deze bachelorproef is om advies te geven aan bedrijven en programmeurs die worstelen met de keuze voor Rexx of Python voor een bepaald probleem. Het is onmogelijk voor elk bestaand probleem een vergelijking te maken. Daarom probeert dit onderzoek enkele veelvoorkomende use cases te gebruiken, die een leidraad kunnen bieden voor de uitwerking van andere use cases.


\section{\IfLanguageName{dutch}{Opzet van deze bachelorproef}{Structure of this bachelor thesis}}%
\label{sec:opzet-bachelorproef}

% Het is gebruikelijk aan het einde van de inleiding een overzicht te
% geven van de opbouw van de rest van de tekst. Deze sectie bevat al een aanzet
% die je kan aanvullen/aanpassen in functie van je eigen tekst.

De rest van deze bachelorproef is als volgt opgebouwd:

In Hoofdstuk~\ref{ch:stand-van-zaken} wordt de huidige stand van zaken beschreven, zodat de lezer een goed idee heeft van de situering van het probleemdomein. Daarnaast wordt er via een literatuurstudie een initiële vergelijking tussen Rexx en Python gemaakt, waar van beide talen de voor- en nadelen worden opgelijst.

In Hoofdstuk~\ref{ch:methodologie} worden de gebruikte onderzoeksmethoden in detail besproken. Die methodologie vormt de basis van het praktijkgedeelte van deze bachelorproef, namelijk de Proof-of-Concept.

% TODO: Vul hier aan voor je eigen hoofstukken, één of twee zinnen per hoofdstuk

In Hoofdstuk~\ref{ch:requirementsanalyse} wordt een requirementsanalyse uitgevoerd. Die zal uitwijzen welke succescriteria zullen worden gebruikt in de Proof-of-Concept om de vergelijking zo accuraat mogelijk te maken. Deze requirementsanalyse bestaat uit literatuurstudie, interviews, en een enquête.

In Hoofdstuk~\ref{ch:Proof-of-Concept} volgt het belangrijkste deel van deze bachelorproef, namelijk de Proof-of-Concept. In de PoC zullen Rexx en Python worden vergeleken aan de hand van enkele veelgebruikte use cases voor systeemprogrammeurs. De use cases en succescriteria worden bepaald met behulp van bevindingen uit de requirementsanalyse.

In Hoofdstuk~\ref{ch:conclusie}, tenslotte, wordt de conclusie gegeven en een antwoord geformuleerd op de onderzoeksvragen. Daarbij wordt ook een aanzet gegeven voor toekomstig onderzoek binnen dit domein.